\begin{frame}{왜? — 앙상블 분산 공식}
  Week 1.1의 \kb{편향-분산 분해}: 예측 오차 = \textbf{편향}(본질적으로 못
  맞히는 부분) + \textbf{분산}(데이터가 조금만 바뀌어도 예측이 흔들리는
  부분). 깊이 제한이 없는 결정 트리는 \textbf{분산이 매우 큼}
  (7.1: \texttt{max\_depth=None} 91.2\% $<$ 깊이 3의 94.7\%).
  \vskip0.8em
  같은 설계의 트리를 $T$개, 각각 다른 부트스트랩 샘플로 학습.
  각 트리의 test 예측 오차 표준편차 $\sigma$, 두 트리 간 상관계수 $\rho$라면
  \textbf{$T$개 평균(다수결 앙상블)의 오차 분산}:
  \[ \mathrm{Var}[\text{앙상블}]
     = \frac{\sigma^2}{T}\,\big(1 + \rho\,(T-1)\big) \]
  (개별 예측 \textbf{평균}의 분산 공식: $T\sigma^2$ + $T(T-1)$개의
  공분산 $\rho\sigma^2$, $T^2$로 나누면 됨.)
\end{frame}
